\section{Conclusión}

El desarrollo de este trabajo práctico permitió analizar las capacidades y limitaciones del paradigma \textit{Greedy} aplicado al problema. 
La implementación de una regla greedy asimétrica, asimétrica porque Sophia maximiza su beneficio local inmediato y minimiza el de Mateo. Demostró ser una solución óptima, logrando una complejidad temporal lineal $\mathcal{O}(n)$ y espacial constante $\mathcal{O}(1)$ que fue corroborada usando el ajuste de Cuadrados Mínimos.

Igualmente, el análisis más relevante de este escenario se basa en comprender las condiciones estrictas que habilitan la optimalidad del enfoque greedy. La estrategia implementada garantiza la victoria absoluta de Sophia únicamente por una restricción crucial de la consigna: ella tiene el control total sobre las decisiones de Mateo, anulando su capacidad de jugar a favor de él. 

En conclusión, el problema propuesto los fundamentos teóricos en el diseño de algoritmos: las soluciones greedy son herramientas ``poderosas'', elegantes y rápidas, pero su éxito está vinculado a la estructura del entorno. 
Una pequeña modificación en las reglas del juego o en el comportamiento de los actores puede llegar a transformar un algoritmo lineal óptimo en una estrategia deficiente por el problema de no tener la visión del ``Big Picture''. 

Se pudo demostrar que la estrategia greedy es un gran diseño y una gran solución, pero también un pensamiento peligroso al ignorar las fronteras del óptimo global.
